% =====================================================================
%  Pacotes e configurações de formatação
% =====================================================================

% ---------- Idioma e codificação ----------
\usepackage[utf8]{inputenc}
\usepackage[T1]{fontenc}
\usepackage[portuguese]{babel}

% ---------- Fontes (Times, como no padrão IEEE) ----------
\usepackage{newtxtext}
\usepackage{newtxmath}
\usepackage{microtype}

% ---------- Matemática e unidades ----------
\usepackage{amsmath}
\usepackage{bm}
\usepackage[
    output-decimal-marker={,},
    group-separator={.},
    per-mode=symbol,
    range-phrase={ a },
    list-final-separator={ e },
    list-pair-separator={ e }
]{siunitx}

% ---------- Figuras, tabelas e diagramas ----------
\usepackage{graphicx}
\usepackage{xcolor}
\usepackage{booktabs}
\usepackage{array}
\usepackage{multirow}
\usepackage{tikz}
\usetikzlibrary{arrows.meta,positioning,calc}

\graphicspath{{./images/}}

% ---------- Listas ----------
\usepackage{enumitem}
\setlist[itemize]{label=\textbullet, leftmargin=*, itemsep=2pt, topsep=4pt}
\setlist[enumerate]{label=\arabic*., leftmargin=*, itemsep=2pt, topsep=4pt}

% ---------- Links e referências cruzadas ----------
\usepackage{url}
\usepackage[autostyle=false,style=english]{csquotes} % aspas “ ” nas referências
\usepackage[hidelinks]{hyperref}

% ---------- Bibliografia (estilo IEEE: [1], [2], ...) ----------
\usepackage[style=ieee,backend=biber,sorting=none,maxnames=6]{biblatex}
\addbibresource{artigo.bib}
\renewcommand*{\bibfont}{\footnotesize}

% ---------- Rótulos em português no padrão IEEE ----------
% Figuras ficam como "Fig. 1." e tabelas como "TABELA I" (padrão do IEEEtran).
\renewcommand{\IEEEkeywordsname}{Palavras-chave}
\addto\captionsportuguese{%
    \renewcommand{\abstractname}{Resumo}%
    \renewcommand{\tablename}{TABELA}%
    \renewcommand{\figurename}{Fig.}%
}




\usepackage{indentfirst}   % recuo também no 1º parágrafo após \section